% Chapter 2 revised: critical literature synthesis, August 2026
\label{ch:theory}

\begin{ataglance}[The argument in outline]
This chapter sets out the theoretical framework. It builds on the literature reviewed in Chapter~\ref{ch:introduction} in seven steps.
\begin{enumerate}\setlength{\itemsep}{2pt}
\item \emph{Build the bridge} from determinants of health to health-related capability, and introduce two illustrative profiles used to check the argument.
\item \emph{Set out the field as a taxonomy.} Five ontological traditions each answer a different question about what ``health'' names (Table~\ref{tab:taxonomy}).
\item \emph{Choose an ontology.} Two candidates, capacity-ceiling and relational, differ on whether context can only convert existing capacity or can also create it.
\item \emph{State the identification problem.} Observed functioning cannot tell a free choice from a blocked one, so capability must be modelled from capacity and context, not read off behaviour.
\item \emph{Place the construct in economics and policy,} where the temptation to overclaim is strongest.
\item \emph{Replace the heuristic formula} with a realisation function derived from stated axioms.
\item \emph{End in falsifiable propositions,} not a declaration of validity.
\end{enumerate}
\end{ataglance}

\noindent
The central problem addressed in this chapter is not whether health is multidimensional. That proposition is well established. Nor is it whether health can be represented by a composite score. The literature contains several successful examples of multidimensional measurement, including the Global Burden of Disease framework, health-expectancy measures, frailty indices, intrinsic-capacity measures, capability measures and latent health metrics developed through factor analysis and item-response theory. The more difficult question is what is gained, and what is lost, when these traditions are brought together in a single construct intended to represent an individual's health-related potential rather than a current health state, a count of deficits, or a probability of an adverse event.

Chapter~\ref{ch:introduction} reviewed five bodies of scholarship: (i) health-state and health-loss measurement; (ii) frailty and deficit accumulation; (iii) the WHO healthy-ageing and intrinsic-capacity framework; (iv) capability theory and capability measurement; and (v) social determinants and social vulnerability. The purpose of bringing these traditions together is not to imply that they are interchangeable. On the contrary, their differences provide the theoretical problem that the thesis must resolve. If intrinsic capacity describes what an individual is able to draw upon, functional ability describes what is realised in context, capability describes the opportunity to achieve valued states and activities, and social vulnerability describes accumulated exposure to social disadvantage, then a proposed health-capability construct must specify precisely where it sits among these concepts.

This leads to a deliberately restrictive claim. The $H_{\text{potential}}$ is not proposed as the first multidimensional health measure, the first latent health score, or the first measure to incorporate social conditions. ATHLOS, for example, constructed a longitudinal health metric using exploratory and confirmatory factor analysis followed by Bayesian multilevel item-response theory, while the WHO framework explicitly links intrinsic capacity, environment and functional ability. \citep{caballero_athlos_2017,who_2015_ageing,who_healthy_ageing_2026} Capability measures such as ICECAP-O likewise demonstrate that capability can be operationalised at the individual level, although the measure is intentionally broader than health. \citep{icecap_o,icecap_o_validation} The proposed contribution is narrower: to test whether a health-bounded construct can explicitly connect biological and functional capacity with enabling social conditions and structural constraints, and whether that construct remains empirically distinguishable from established measures of intrinsic capacity, frailty, social vulnerability and current health.

The chapter consequently ends not with a declaration that the proposed construct is valid, but with a set of propositions that the empirical work must attempt to falsify. The logic is important. A thesis of this kind is strongest when the data are allowed to demonstrate that the proposed construct is unnecessary, multidimensional without a valid general factor, or prognostically redundant. The objective is not to protect the proposed index from such findings but to determine what the evidence can legitimately support.

\section{From determinants to capability: the proposed bridge}
\label{sec:intro-bridge}


The central conceptual move of the thesis is to distinguish between a determinant of health and a component of health-related capability. Income, education, housing, transport and access to care can influence health without necessarily being equivalent to health. Conversely, intrinsic capacity may be high while opportunities to exercise that capacity are constrained by an inaccessible environment. The same underlying physical capacity may therefore produce very different lived possibilities in different contexts.

This distinction is consistent with the capability approach, but the thesis does not equate health capability with the broader philosophical concept of human wellbeing. Capability theory is concerned with substantive opportunities, whereas health measurement requires a bounded account of what those opportunities mean in relation to health. The WHO framework offers a practical bridge because it explicitly describes functional ability as emerging from intrinsic capacity, environmental characteristics and their interaction \citep{who_2015_ageing}. The proposed construct builds from this relationship but asks whether it can be made sufficiently explicit and measurable to support individual-level statistical inference.

A second bridge is required with social-determinants and life-course research. Social disadvantage is not simply an additional ``risk factor'' to be placed beside age, multimorbidity or cognition. Education, material resources, housing, social relationships and access to services are embedded in cumulative life-course processes and in institutions that distribute opportunities unequally \citep{who_sdh_2008,marmot_2005}. If these conditions are incorporated into a health-capability measure, the thesis must therefore avoid implying that structural disadvantage is a personal deficit. The proposed interpretation is instead that constrained circumstances can reduce the extent to which existing health capacity can be realised.

A third bridge concerns frailty and reserve. Frailty research provides a powerful demonstration that health can be represented as multidimensional vulnerability rather than as a single diagnosis. At the same time, frailty and capability answer different questions. Frailty asks, in broad terms, how vulnerable a person is to adverse stressors; capability asks what the person is able to realise within a given context. The empirical relationship between these constructs is therefore a testable question, not a definitional identity. A successful H-potential measure should correlate with frailty while retaining information that is not reducible to accumulated deficit burden.

These bridges lead to the proposed construct, $H_{\text{potential}}$, introduced in Chapter~\ref{ch:introduction}. It is defined as health-related capability in later life: the combination of biological and functional capacities available to an individual and the enabling or constraining conditions through which those capacities can be realised. The definition is deliberately relational. It does not treat social disadvantage as a biological property, nor does it treat future mortality as part of the definition of health.

Before the theory is developed, two illustrative profiles are introduced against which each later claim can be checked.

\section{Two illustrative profiles}
\label{sec:intro-cases}

Abstract argument about health measurement is easy to agree with and hard to
test. To keep the argument honest, two people are carried through the whole
thesis. Neither is a real patient or a case study; both are composites assembled from the
distributions in the data sources used here. They are introduced now, applied
to every equation as it appears, and revisited in the final chapter.

\begin{ataglance}[Mrs~A, the population case]
Mrs~A is 78, widowed, and lives alone on the second floor of a walk-up flat.
She has knee osteoarthritis, mild cognitive impairment, hypertension and early
chronic kidney disease. She dresses and washes herself but no longer carries
shopping up the stairs. Her income is a state pension. She does not drive; her
clinic is forty minutes away by two buses, and she has missed three of her last
eight appointments. Her daughter telephones daily and visits monthly from
another city.

Mrs~A stands for the population the empirical work actually studies: an older
adult of the kind sampled by the National Health and Aging Trends Study, whose
clinical problems are unremarkable and whose circumstances do most of the work.
She is the case used whenever the question is measurement, prediction or
population comparison.
\end{ataglance}

\begin{ataglance}[Mr~B, the translational case]
Mr~B is 58 and has severe haemophilia~A with a high-titre inhibitor. He has a
target right ankle and a destroyed left knee from bleeds sustained before
prophylaxis was available to him, chronic pain, and the metabolic consequences
of two decades of reduced mobility. He is employed part-time. Whether he
receives emicizumab or on-demand bypassing agents depends on which country he
lives in and which formulary applies to him.

Mr~B stands for the translational question: what happens to a construct built
in one population when it is carried to a disease, a health system and a
funding regime it was not developed in. He is the case used whenever the
question is transportability, treatment effect or the price of a decision.
\end{ataglance}

The two are chosen because they fail in different directions. Mrs~A has modest
biology and hostile circumstances. Mr~B has severe, well-characterised biology
and circumstances that swing between generous and punitive depending on
jurisdiction. A construct that describes only one of them well has not done the
job this thesis sets out.

\section{What Does It Mean to Measure Health?}

Measurement begins with a construct, not an algorithm. A score can be calculated whenever observations are combined, but calculation does not establish that the resulting quantity corresponds to a meaningful latent attribute. This distinction is particularly important in multidimensional health research because the word \emph{health} is used to describe states, capacities, functionings, opportunities, vulnerabilities and outcomes that occupy different conceptual levels.

The WHO healthy-ageing framework provides an important starting point because it explicitly distinguishes intrinsic capacity from functional ability. Intrinsic capacity refers to the composite of an individual's physical and mental capacities, while functional ability emerges from the interaction between those capacities and relevant environmental characteristics. The environment is defined broadly, including physical, social and policy contexts and the systems and relationships through which people live. \citep{who_2015_ageing,who_healthy_ageing_2026} This formulation prevents a common measurement error: interpreting an observed limitation as evidence that the underlying biological capacity is absent when the observed limitation may instead reflect an environmental barrier.

The distinction can be expressed conceptually, for person $i$, as
\begin{equation}
\boxed{FA_i=f(IC_i,E_i,IC_i\times E_i)},
\label{eq:who_fa_conceptual}
\end{equation}
where $FA_i$ denotes realised functional ability, $IC_i$ intrinsic capacity and
$E_i$ the relevant environment. Equation~\ref{eq:who_fa_conceptual} is a
translation of the WHO conceptual architecture into notation; it is not an
equation estimated or endorsed by WHO. Writing the interaction among the
arguments of $f$ does not assert a linear interaction term, a multiplicative
production function or any particular scale.

\begin{definition}[Functional ability in the WHO framework]
\label{def:lit-fa}
For person $i$, let $IC_i$ denote intrinsic capacity, $E_i$ the relevant
environment and $FA_i$ realised functional ability. The WHO healthy-ageing
architecture, written in the notation of \eqref{eq:who_fa_conceptual}, posits a
function $f$ with
\[
  FA_i=f\bigl(IC_i,E_i,IC_i\times E_i\bigr).
\]
The definition fixes no scale for $IC$ or $E$, no sign for the interaction and
no functional form for $f$. \citep{who_2015_ageing,who_healthy_ageing_2026}
\end{definition}

\begin{assumption}[The interaction is not redundant]
\label{ass:lit-interaction}
In Definition~\ref{def:lit-fa}, $f$ is not additively separable in its first
two arguments: there exist profiles at which
$\partial^2 f/\partial IC\,\partial E\neq0$.
\end{assumption}

Stated as an assumption rather than folded into the definition, because it is
the substantive claim and it can fail. If the cross-partial were zero
everywhere, capacity and environment would contribute independently, the third
argument of $f$ would carry no information, and an additive index would lose
nothing. Everything the thesis goes on to build rests on this assumption being
false, and Chapter~\ref{ch:results2} is where it is put at risk.

\begin{reading}[Why the interaction term is the whole argument]
Strip \eqref{eq:who_fa_conceptual} of the notation and it says three things.
What a person can do depends on their body. It depends on their surroundings.
And, third, it depends on the two together in a way neither explains alone.
The first two are uncontroversial. The third is the claim.

The clinical case for it is easy to state. Two women have identical hip
radiographs, identical grip strength and identical gait speed: the same $IC$.
One lives in a bungalow with a stairlift-free front door and a supermarket
across the road. The other lives on the second floor of a walk-up. Their $FA$
differs, and it differs by more than either the hips or the stairs would
predict separately, because the stairs matter more when the hips are worse.
That excess is $IC_i\times E_i$.

A measure that omits the interaction is committed to saying the stairs cost
every woman the same. They do not, and the cost of pretending they do falls
hardest on the women for whom the answer matters most.
\end{reading}

The conceptual sequence used in this thesis is \citep{who_2015_ageing}
\begin{equation}
\boxed{\begin{aligned}
\text{health conditions}&\longrightarrow\text{intrinsic capacities}
\longrightarrow\text{health-related capabilities}\\
&\longrightarrow\text{functional ability}
\longrightarrow\text{observed outcomes}
\end{aligned}},
\label{eq:conceptual_sequence}
\end{equation}
with environmental and social structures acting across the sequence. The
arrows express theoretical ordering, not identified causal effects.

\begin{definition}[The conceptual sequence]
\label{def:lit-sequence}
The sequence of \eqref{eq:conceptual_sequence} is the ordered chain
\[
  \text{health conditions}\rightarrow IC\rightarrow
  \text{health-related capability}\rightarrow FA\rightarrow
  \text{observed outcomes},
\]
with environmental and social structure acting on every arrow. Each arrow is a
many-to-one map, so information present at one stage need not survive to the
next. The ordering is theoretical: it says which quantity is upstream of which,
and it does not assert an identified causal effect at any arrow.
\end{definition}

Each arrow in \eqref{eq:conceptual_sequence} is a place where information is
lost, and naming the losses is the point of writing the sequence out. Mrs~A's
osteoarthritis is a health condition. Her reduced quadriceps strength and
slowed gait are the intrinsic capacities it produces. Her capability to leave
the flat is what those capacities plus a second-floor walk-up amount to. Her
functional ability is what she in fact does, which is to go out twice a week.
Her observed outcome is a missed clinic appointment recorded in an
administrative dataset.

Reading backwards along that chain is where measurement goes wrong. The missed
appointment is compatible with severe disease, with a difficult journey, or
with a woman who had something better to do. Existing measures mostly sample
the right-hand end of the sequence and infer leftwards. This thesis samples the
left and middle and refuses to infer from behaviour, which is the position the
next equation formalises.

The thesis therefore asks whether underlying health-related potential can be
measured separately from realised function. Its general construct statement is
\begin{equation}
\boxed{H_{P,i}=f(BP_i,SM_i,SB_i)},
\label{eq:hp_general_class}
\end{equation}
where $BP_i$ is biological and functional capacity, $SM_i$ comprises social and
environmental enabling conditions, and $SB_i$ comprises structural and social
barriers. The mapping $IC\rightarrow BP$ is approximate rather than an asserted
equivalence. Likewise, $E\rightarrow(SM,SB)$ separates enabling conditions from
constraints without implying that every environmental variable belongs inside
the construct. Whether social structure is a component, determinant or modifier
is tested through competing specifications.

\begin{definition}[The H-potential construct class]
\label{def:lit-hpclass}
An H-potential representation is any function $f$ of three domain scores,
\[
  H_{P,i}=f\bigl(BP_i,SM_i,SB_i\bigr),
\]
where $BP_i$ is biological and functional capacity, $SM_i$ enabling social and
environmental conditions and $SB_i$ structural and social barriers, as in
\eqref{eq:hp_general_class}. The definition names a class, not a member. It
commits to two claims and no more: that the three domains partition what
matters for health potential, and that they can be measured apart. Choice of
$f$ within the class is the subject of Section~\ref{sec:lit-candidates} and
Chapter~\ref{app:crf}.
\end{definition}

\begin{reading}[What \eqref{eq:hp_general_class} does and does not commit to]
The equation is deliberately almost empty. It names three domains and asserts
that health potential is some function of them. It does not say the function is
addition, or multiplication, or anything else; $f$ is unspecified on purpose,
because choosing its form is a separate argument made in
Chapter~\ref{app:crf} and not smuggled in here.

Two things are being claimed even so, and both can fail. The first is that
these three domains are the right partition: that nothing important about
health potential falls outside biological capacity, social enablement and
structural burden. The second is that they are separable enough to be measured
apart. If the NHATS indicators cannot distinguish social enablement from
structural burden, the partition collapses to two domains and the thesis has to
say so.

Note also what sits outside the equation. There is no time index, so this is a
snapshot; the dynamic version is Appendix~\ref{app:dynamical}. And there is no
outcome on either side, which is the separation the whole design turns on.
\end{reading}

\begin{table}[htbp]
\centering
\caption{Conceptual and empirical roles in the H-Potential argument}
\label{tab:concept_roles}
\small
\begin{tabular}{p{0.22\linewidth}p{0.22\linewidth}p{0.44\linewidth}}
\hline
Concept & Mathematical representation & Empirical role \\
\hline
Intrinsic capacity & $IC$ & Latent or domain-specific capacities; principal clinical predecessor \\
Environment & $E$ & Context, enablement or constraint; not presumed reflective \\
Functional ability & $FA=f(IC,E,IC\times E)$ & Realised capability and functioning \\
Frailty & $F$ & Competing formative/composite construct \\
Social vulnerability & $SV$ & Contextual competitor and possible determinant \\
H-Potential & $H_P=f(BP,SM,SB)$ & Proposed health-bounded construct \\
Mortality & $(T,D)$ & Downstream prognostic outcome, never a defining indicator \\
Hospitalisation & $(T_H,D_H)$ & Downstream outcome with death as a competing event \\
QALY & utility $\times$ survival & Cardinal economic comparator, not a latent-score interpretation \citep{whitehead_2010,drummond2015} \\
Health expectancy & survival/function & Population-level comparator \\
\hline
\end{tabular}
\end{table}

The capability approach provides a second conceptual bridge. Sen distinguishes achieved functionings from the capability to achieve alternative functionings. \citep{sen_1993} This distinction is useful in health because the same observed outcome can arise from different opportunity structures. A person may not participate in an activity because of choice, because of impaired capacity, or because environmental and institutional constraints make participation difficult. A measurement system that observes only achieved functioning can therefore conflate inability with constrained opportunity.

However, the capability approach cannot be imported into health measurement without qualification. Capability is a theory of human freedom and wellbeing, not a ready-made psychometric definition of health. ICECAP-O illustrates this boundary explicitly: it was developed as a measure of capability wellbeing in older people and intentionally extends beyond health states. \citep{icecap_o,icecap_o_validation} Empirical comparisons also suggest that ICECAP-O and direct health measures overlap only partially, particularly in physical functioning. \citep{icecap_o_validation} This is not a weakness of ICECAP-O; it is evidence that capability and health are related constructs with different scopes.

The present thesis therefore uses the term \emph{health-related capability} in a deliberately restricted sense. It refers to the capacity and opportunity to sustain health-related functioning under the biological and social conditions in which an individual lives. It is not intended to measure total human capability, flourishing, life satisfaction or general wellbeing. This boundary is necessary for discriminant validity: if the construct can only be distinguished from general wellbeing by its name, its theoretical specification is inadequate.

A further distinction is required between a latent construct and a prognostic outcome. A score derived from NHATS is not itself future health. It is an estimate of a latent state hypothesised to be associated with subsequent outcomes. Mortality and hospitalisation therefore cannot be used to define the construct without creating circularity. The measurement model must be developed before the prognostic model is evaluated. Prediction is evidence relevant to validity, but prediction alone does not establish construct meaning.

If several conceptions of health are in play, the first task is to lay them out and see which question each answers. The next section does that side by side rather than in sequence, because the differences between them are the argument.

\section{A Taxonomy of Ontological Traditions}

Five ontological traditions define what the word \emph{health} names, and each answers
a different question. To make the gap the thesis addresses precise, they are set out as
a taxonomy rather than a list, alongside the health-economic, social-vulnerability and
capability-wellbeing measures that serve as comparators. Table~\ref{tab:taxonomy}
states, for each, the construct it targets, how it is structurally represented, its
conceptual limit, and the role it plays in the H-potential argument. The point of the
taxonomy is negative: while every entry is well founded for its own question, none
isolates the object the thesis tests, namely health-related capability as biological
capacity converted, or blocked, by enabling and constraining context.

\begin{table}[htbp]
\centering
\caption{A taxonomy of ontological traditions in health measurement, and the role of each in the H-potential argument.}
\label{tab:taxonomy}
\footnotesize
\renewcommand{\arraystretch}{1.15}
\begin{tabular}{p{0.17\linewidth} p{0.17\linewidth} p{0.17\linewidth} p{0.20\linewidth} p{0.20\linewidth}}
\toprule
\textbf{Tradition} & \textbf{Primary construct target} & \textbf{Structural representation} & \textbf{Conceptual limit} & \textbf{Role in the argument} \\
\midrule
Biostatistical naturalism (Boorse \citep{boorse_1977}) & Statistically typical biological function vs a reference class & Species-typical functional norms & Places social context outside health & Fixes the biological pole; motivates $BP$ but is too narrow alone \\
\addlinespace
Action-theoretic holism (Nordenfelt) & Ability to reach vital goals in standard circumstances & Agent-and-goal ability & ``Standard circumstances'' hides structural oppression & Introduces agency; shows why context must be explicit \\
\addlinespace
Capability approach (Sen; Venkatapuram) & Substantive opportunity to achieve valued functionings & Functionings vs the capability set & Broader than health; risks construct conflation & Supplies the capability/functioning distinction, bounded here to health \\
\addlinespace
Deficit accumulation / frailty (Rockwood \citep{rockwood_mitnitski_2007,searle_2008}; Mitnitski) & Accumulated physiological vulnerability & Formative ratio $FI=\sum\text{deficits}/\text{total}$ & Agnostic on cause; omits enabling structures & The competing composite for risk stratification \\
\addlinespace
Intrinsic capacity / healthy ageing (WHO) & Internal physical and mental reserve & Reflective/formative domains & Heterogeneous operationalisation; excludes external context & The biological foundation, $BP$ \\
\midrule
Health-economic summaries (QALY / DALY) & Evaluative preference utility, or disease loss & Composite ratio/interval scale & Evaluates consequences; omits enabling context & Cardinal economic comparator \\
\addlinespace
Social vulnerability & Contextual disadvantage & Composite index of social deficits & Treats context as a personal deficit & Direct contextual competitor \\
\addlinespace
Capability wellbeing (ICECAP-O) \citep{icecap_o,coast_2008} & General human flourishing & Preference-weighted index & Exceeds the health domain & Boundary marker for health-specific freedom \\
\midrule
\textbf{$H_{\text{potential}}$} & Health-related capability in context & Axiomatic non-compensatory realisation function & Requires multi-stage identification and external calibration & The proposed construct \\
\bottomrule
\end{tabular}
\end{table}

\begin{reading}[How to read the taxonomy]
Each row is legitimate for the question it was built to answer; the table is not a
list of failures. Its work is to locate a gap: read down the ``conceptual limit''
column and every established tradition either excludes context, or dissolves the
boundary of health, or treats disadvantage as a personal deficit. The final row is
the thesis's wager that a health-bounded, non-compensatory construct can sit in that
gap, and the rest of the chapter tests whether it can be told apart from its
neighbours.
\end{reading}

A gap in a table is not yet a justification. Whether anything belongs in that
gap depends on what one thinks health is, and on that question the philosophy
of medicine is openly divided.

\section{Rival Philosophies of Health and Their Mathematical Consequences}

The word \emph{health} does not name a philosophically uncontested natural
kind. A naturalist position, associated most prominently with Boorse's
biostatistical theory, locates health in statistically typical biological
functioning relative to an appropriate reference class. Contemporary defences
retain the attraction of separating biological description from judgements
about whether a condition is undesirable. \citep{schwartz_2014_boorse} On this
view, housing, transport, income and institutional access influence welfare and
the consequences of disease, but they are not themselves parts of health.

This position supplies an important objection to H-Potential. If health is
biological functional efficiency, placing social enablement inside a health
score commits a category error. The naturalist model would instead write
\[
H_i^{N}=P_i,
\qquad
FA_i=f(P_i,E_i),
\]
so that $P_i$ measures health and $E_i$ helps explain realised functional
ability. This is a serious comparator, not a straw position. It protects the
medical construct from expanding until every social good becomes health.

Applied to Mr~B, the naturalist reads his health off his factor level, his
joint scores and his inhibitor titre, and says his formulary is a fact about
his country rather than a fact about him. There is something right in that. A
haematologist who recorded ``poor drug funding'' in the problem list alongside
``high-titre inhibitor'' would be making a category error, and a health measure
that let policy failures show up as sickness in the patient would be doing the
same thing at scale. The objection has to be answered rather than waved
through, and the answer this thesis gives is the layered position at the end of
this section: report $P$ separately, always.

Nordenfelt's action-theoretic account moves in a different direction. Health is
connected to a person's ability, under accepted circumstances, to achieve vital
goals. \citep{nordenfelt_1995} This captures something the biological account
misses: intact organs do not guarantee the practical ability to act. Yet it
introduces difficult questions. Which goals are vital? Which circumstances
count as accepted or standard? If an inaccessible environment is treated as
standard, injustice may be normalized; if every adverse circumstance is
incorporated, health risks becoming indistinguishable from general welfare.

Sen's capability approach separates achieved functionings from the real
opportunities to achieve them. \citep{sen_1993} It therefore explains why the
same observed functioning may arise from choice, biological limitation or
constraint. Venkatapuram extends this reasoning by presenting health capability
as a meta-capability shaped by biomedical and social causes and connected to
human dignity and justice. \citep{health_capability_right,venkatapuram_2011}
This supports a relational model in which capability is jointly constituted:
\[
H_i^{R}=G(P_i,S_i,R_i),
\]
without assuming that biomedical capacity is an external ceiling.

The clinical intuition behind the relational reading is also familiar. A person
with paraplegia and a well-fitted wheelchair in an accessible city is not
described by their spinal cord alone; a person with the same lesion, no
wheelchair and no ramps is differently placed in a way that ``same disease,
different circumstances'' does not quite capture. The relational account says
the difference is in their health capability, not merely alongside it.

Daniels provides a further bridge to policy. On a Rawlsian account, protecting
health matters because departures from normal functioning restrict the range of
opportunities open to people; fair health policy must therefore address both
health needs and fair procedures when not all needs can be met.
\citep{daniels_2007} This approach does not imply that one composite score can
settle priority. It makes procedural legitimacy, reasons, revisability and
opportunity effects part of the assessment.

These positions generate two rival ontologies for the thesis. The
\emph{capacity-ceiling ontology} treats $P$ as the maximum potentially
realisable quantity and context as a conversion factor:
\[
H^{C}(P,S,R)=P\,M(S,R),\qquad 0\le H^{C}\le P.
\]
The \emph{relational ontology} permits technology, assistance and institutions
to enlarge effective capability beyond what an unaided capacity measure would
predict:
\[
H^{R}(P,S,R)=G(P,S,R),\qquad H^{R}\not\le P\text{ necessarily}.
\]
Neither is selected by terminology. The ceiling model is more conservative and
clinically anchored; the relational model is more faithful to capability and
disability scholarship. The empirical analysis can examine their consequences,
but fit cannot decide the philosophical question by itself.

\begin{definition}[Capacity-ceiling ontology]
\label{def:lit-ceiling}
Write the profile as $(P,S,R)$ with $R=1-B$. The capacity-ceiling ontology
takes biomedical capacity to be the maximum realisable quantity and context to
be a conversion factor $M:[0,1]^2\rightarrow[0,1]$:
\[
  H^{C}(P,S,R)=P\,M(S,R),\qquad 0\le H^{C}(P,S,R)\le P .
\]
The upper bound is the substantive content. Context can convert capacity that
exists and cannot create capacity that does not, so $P=0$ forces $H^{C}=0$ for
every $S$ and $R$.
\end{definition}

\begin{definition}[Relational ontology]
\label{def:lit-relational}
The relational ontology drops the ceiling. Capability is jointly constituted by
capacity and context through a function $G$,
\[
  H^{R}(P,S,R)=G(P,S,R),
\]
with no requirement that $H^{R}\le P$. Technology, assistance and institutions
may therefore enlarge effective capability past what an unaided capacity
measure predicts, and $P=0$ need not force $H^{R}=0$.
\end{definition}

\begin{reading}[What the two definitions disagree about]
Definitions \ref{def:lit-ceiling} and \ref{def:lit-relational} differ in one
inequality. The ceiling says a wheelchair cannot give a person legs; it can
only let the legs they have carry them further. The relational account says
that a person with a chair, a ramp and a lift is able to do things, and that
being able to do things is what health capability was supposed to mean.

Both are defensible. Neither is chosen here. The thesis carries both, reports
where they diverge, and treats the divergence as a result rather than a
nuisance.
\end{reading}

\begin{reading}[The one case that separates them]
The two ontologies agree almost everywhere and disagree at a boundary. Take a
person with $P$ at or near zero and excellent circumstances. Under the ceiling
model, $H^C=P\,M(S,R)$, and multiplying zero by anything gives zero: their
health capability is nil regardless of the ramps. Under the relational model,
$G$ need not vanish, so their circumstances still count.

Which answer is right is not a question data can settle, because the two
models make identical predictions everywhere else and this corner of the
profile space is sparsely observed. Choosing by fit would mean letting the
distribution of a convenience sample decide a question in the philosophy of
medicine, which is why the thesis carries both and reports where they diverge
rather than picking one and calling the choice a finding.
\end{reading}

The thesis therefore adopts a layered position. Biomedical Potential is
reported as a distinct clinical domain. The multidomain profile is primary.
Scalar aggregation is secondary and conditional. Both ceiling and relational
aggregations are sensitivity models. This prevents a philosophical commitment
from being concealed inside a single equation.

Both ontologies share one difficulty, and it is prior to either of them. Each
treats capability as something a person has. Nobody has ever observed one.

\section{Epistemology: Why Capability Cannot Be Read Directly from Function}

Let $\mathcal C_i$ be the set of health-related functionings genuinely
available to person $i$, and let the achieved functioning be
\[
F_i=\chi_i(\mathcal C_i),
\]
where $\chi_i$ represents choice under the person's information, preferences
and circumstances. It follows that
\[
F_i=F_j\not\Rightarrow\mathcal C_i=\mathcal C_j.
\]
Two people may display the same activity while one had several genuine options
and the other had none. Conversely, non-achievement does not establish
incapability.

\begin{definition}[Achieved functioning map]
\label{def:lit-choicemap}
Let $\mathcal U$ be the set of health-related functionings and let
$\mathcal C_i\subseteq\mathcal U$ be the capability set of person $i$, meaning
the functionings genuinely available to them. The achieved functioning is
\[
  F_i=\chi_i(\mathcal C_i),\qquad \chi_i(\mathcal C)\in\mathcal C ,
\]
where the selector $\chi_i$ represents choice under the person's information,
preferences and circumstances. Only $F_i$ is recorded.
\end{definition}

\begin{theorem}[Observed functioning does not identify the capability set]
\label{thm:lit-identification}
Let $\mathcal C$ range over the non-empty subsets of $\mathcal U$ and let
$\chi$ be a selector in the sense of Definition~\ref{def:lit-choicemap}. If
$|\mathcal U|\ge2$ then $\chi$ is not injective, and consequently
\[
  F_i=F_j\;\not\Rightarrow\;\mathcal C_i=\mathcal C_j .
\]
No function of observed functioning alone recovers $\mathcal C_i$. For any
candidate inverse $g$ with $g(F)\in2^{\mathcal U}$, and any $F$ chosen from $k$
distinct capability sets, $g$ returns the right set for at most one of the $k$
people and is wrong for the other $k-1$.
\end{theorem}

\begin{proof}[Proof sketch]
Take $u\in\mathcal U$ and any two distinct sets $\mathcal C_i,\mathcal C_j$ that
both contain $u$; a selector with $\chi_i(\mathcal C_i)=\chi_j(\mathcal
C_j)=u$ exists whenever $|\mathcal U|\ge2$, so $\chi$ is not injective and the
implication fails. The counting claim follows because $g$ is a function: it
assigns one value to the argument $F$, and that value equals at most one of the
$k$ distinct sets. Identification therefore requires information outside $F$,
which is what Definition~\ref{def:lit-hpclass} supplies by modelling capability
from capacity and context.
\end{proof}

\begin{reading}[What the theorem does not say]
It does not say observed functioning is uninformative. A record of what someone
did is evidence about what they could have done, and in a large sample it
carries a great deal. The theorem says something narrower and harder: the
evidence does not determine the answer, so the gap has to be closed by
assumption, and the assumption has to be stated.

It also sets the price of the alternative. Modelling capability from capacity
and context replaces an unidentified quantity with an estimated one. That is
progress only if the model is written down, published and allowed to fail.
\end{reading}

\begin{reading}[The identification problem, at the bedside]
This is the most consequential piece of algebra in the chapter and it says
something a clinician already knows. Observed behaviour is a choice made from
an opportunity set, and only the choice is recorded.

Three people did not attend the diabetic eye-screening appointment. The first
could not afford the bus. The second was too breathless to travel. The third
read the invitation, weighed it, and decided her cataract surgery next month
made it pointless. The administrative dataset records one variable, DNA, and
$F_i$ is identical for all three. Their opportunity sets $\mathcal C_i$ are
completely different, and so is what should be done about them.

Now run it the other way. Two men both walk four hundred metres. One does it
easily. One does it because he must collect his own medication and has spent
the rest of the day recovering. Same $F$, different $\mathcal C$. Every measure
built on observed functioning, which includes most activities-of-daily-living
scales, is blind to this by construction.

The formal statement $F_i=F_j\not\Rightarrow\mathcal C_i=\mathcal C_j$ is the
reason the thesis models capability from capacity and context rather than
reading it off behaviour. It also states the price: capability is not directly
observable, so it must be inferred, and an inference can be wrong in ways an
observation cannot.
\end{reading}

\begin{workedexample}[lit-identification]{Three people, one recorded outcome}
Three women were invited to diabetic eye screening and none attended. The
dataset records the same value of $F$ for all three: did not attend. Their
capability sets are not the same.

\begin{center}\small
\begin{tabular}{p{0.20\linewidth}p{0.40\linewidth}p{0.28\linewidth}}
\toprule
Person & Capability set $\mathcal C_i$ & Recorded $F_i$ \\
\midrule
No fare & \{home\} & did not attend \\
Breathless & \{home, attend\} & did not attend \\
Informed refusal & \{home, attend, attend and shop\} & did not attend \\
\bottomrule
\end{tabular}
\end{center}

\noindent The sets have one, two and three elements. Theorem
\ref{thm:lit-identification} applies with $k=3$: any rule that reads a
capability set off the record returns one set, so it is right for one woman and
wrong for two. The clinical responses differ completely. The first needs the
fare, the second needs transport and a review of her breathlessness, the third
needs nothing at all.

\begin{wecode}[Python]
sets = {
    "no fare":          frozenset({"home"}),
    "breathless":       frozenset({"home", "attend"}),
    "informed refusal": frozenset({"home", "attend", "attend and shop"}),
}
observed = {name: "home" for name in sets}       # chi_i(C_i), identical
print(len(set(observed.values())), "recorded value(s)")
print(len(set(sets.values())), "distinct capability sets")
for guess in sets.values():
    right = sum(1 for c in sets.values() if c == guess)
    print("a fitted inverse gets", right, "right and", len(sets) - right, "wrong")
\end{wecode}

\weanswer{One recorded value, three distinct capability sets, and any inverse
map is wrong for two of the three. Asserted by
\texttt{test\_theorem\_lit\_choice\_map\_is\_not\_injective}.}
\end{workedexample}

This is an identification problem. Cross-sectional functioning cannot uniquely
recover an unobserved opportunity set. Self-report does not solve it completely
because expectations can adapt to long-standing deprivation, while external
assessment may substitute professional values for those of the person. The
claim made by H-Potential must therefore remain modest: it estimates selected
conditions relevant to health-related opportunity; it does not directly observe
freedom.

The implication for validity is that statistical covariance is insufficient.
Content validation must include clinicians, patients and affected communities;
the indicators must be checked against what people have reason to value; and
discordance between observed functioning, reported opportunity and structural
conditions must be studied rather than averaged away.

With the conceptual and epistemic ground laid, and the measurement traditions
reviewed in Chapter~\ref{ch:introduction}, the question becomes what such a
measure could legitimately be used for. The economics is where the temptation
to overclaim is strongest.

\section{Economic Interpretation: Value, Opportunity Cost and Distribution}

H-Potential is not an economic evaluation outcome. A QALY combines time and
preference-based health-state value; a budget-impact analysis measures
affordability; cost-effectiveness compares incremental costs and consequences.
An H-Potential score has none of these meanings by construction. In particular,
the latent scale does not support a cost-per-H-Potential-unit threshold.

The economic relevance instead lies in diagnosis and heterogeneity. Suppose an
intervention $a$ changes the profile from $X_i(0)$ to $X_i(a)$ at cost $C_i(a)$.
Its profile effect is
\[
\Delta X_i(a)=X_i(a)-X_i(0),
\]
and a scalar effect $\Delta H_i(a)$ is calculated only under a declared
aggregator. Economic evaluation must then compare intervention costs and
established consequences, while asking who gains and which constraint is
changed. A treatment may raise $P$; accessible transport may raise $S$; removal
of financial barriers may raise $R$. Equal total effects can therefore imply
different budgets, responsible agencies and policy mechanisms.

\begin{reading}[Why the profile effect is kept as a vector]
Three interventions produce the same $\Delta H$ of $+0.05$ for Mrs~A. A knee
replacement raises $P$. A befriending scheme raises $S$. Waiving her clinic
co-payment raises $R$. A scalar records all three as identical and hands the
decision-maker no information at all.

The vector $\Delta X$ records them as $(+0.05,0,0)$, $(0,+0.05,0)$ and
$(0,0,+0.05)$, and the three go to different budgets. The first is an
orthopaedic list with an eighteen-month wait. The second is a local-authority
grant of a few hundred pounds. The third is a line in a finance director's
spreadsheet and costs almost nothing to change. Equal capability effects,
utterly unequal in cost, in who is responsible, and in how long they take.

This is the practical reason the framework refuses to collapse the profile.
Not because a scalar is philosophically impure, but because the scalar deletes
the only column a commissioner can act on.
\end{reading}

Opportunity cost remains unavoidable. Funding an intervention displaces other
health-producing activity. H-Potential cannot make that displaced health
disappear and must not be used as a moral multiplier attached after a
cost-effectiveness result. If allocation is considered, the minimum economic
record includes incremental cost, QALYs or another justified consequence,
budget impact \citep{sullivan_2014}, displaced activity, uncertainty and the distribution of gains
and losses.

Distributional cost-effectiveness analysis provides the closest established
economic bridge. It extends conventional evaluation by estimating how health
and opportunity costs are distributed across equity-relevant groups and by
making inequality judgements explicit. \citep{asaria_2016_dcea} H-Potential
profiles could help describe baseline constraints or intervention mechanisms,
but should not replace DCEA's health-opportunity-cost accounting. Nor should a
low score automatically generate priority: doing so could favour interventions
with little benefit, penalise those whose capability is difficult to measure or
double-count disadvantage already included in an equity analysis.

Three distinct economic uses must therefore be separated:
\begin{enumerate}
 \item \emph{descriptive use}: identify the domain in which capability is
 constrained;
 \item \emph{stratified evaluation}: estimate costs and outcomes separately
 for relevant profile groups;
 \item \emph{allocative use}: apply explicit social value judgements through
 a legitimate appraisal process.
\end{enumerate}
Only the first is supported by construct validation alone.

The distinction is worth stating bluntly, because the slide from the first use
to the third is how measures of this kind cause harm. Establishing that Mrs~A's
constraint sits in the structural domain is a descriptive finding and the
validation work in this thesis could support it. Concluding that she should
therefore receive funding ahead of someone else is an allocative judgement, and
nothing in a factor analysis licenses it. The sections that follow set out what
would have to be true, and what procedures would have to be in place, before
the framework could be used at the third level at all.

\subsection{The policy evidence vector and its calculations}

Because the foregoing distinctions are easy to lose once a score enters a
decision process, the minimum policy object can be written explicitly. Let
$a$ denote an intervention and let $j$ denote the jurisdiction in which it is
being considered. The evidence required for appraisal is the vector
\[
\mathcal D_j(a)=
\left(
\Delta X_j(a),\Delta Q_j(a),\Delta C_j(a),
\operatorname{BI}_j(a),\operatorname{Dist}_j(a),
\operatorname{Unc}_j(a),\operatorname{Gov}_j(a)
\right).
\]
This notation is not intended to create a new omnibus score. Its purpose is the
opposite: it prevents analytically different considerations from disappearing
inside a single number. The coordinates record, respectively, change in the
H-Potential domain profile, change in an established health outcome,
incremental cost, budget impact, distributional consequences, uncertainty and
governance requirements. This follows the established distinction between
economic evaluation and affordability analysis and the multidisciplinary logic
of HTA. \citep{drummond2015,who_hta_2014,who_hta_institutional_2021}

The first coordinate is genuinely multidimensional. With
\[
X_j(a)=\bigl(P_j(a),S_j(a),R_j(a)\bigr),
\]
where $P$ denotes biomedical and functional capacity, $S$ social and
environmental enablement, and $R=1-B$ freedom from structural burden, the
profile change is
\[
\Delta X_j(a)=X_j(a)-X_j(0)
=\bigl(\Delta P_j(a),\Delta S_j(a),\Delta R_j(a)\bigr).
\]
Larger coordinates are favourable. The vector should normally be reported in
this disaggregated form. A scalar $\Delta H_j(a)$ may be added only when the
aggregation rule, weights and compensation assumptions have been defended; it
does not follow from subtraction alone.

The economic coordinates use conventional incremental analysis. For a time
horizon $T$, discount factor $d_t=(1+r)^{-t}$, intervention arm $a$ and
comparator $0$,
\[
\Delta Q_j(a)=\sum_{t=0}^{T}d_t
  \{Q_{jt}(a)-Q_{jt}(0)\},
\qquad
\Delta C_j(a)=\sum_{t=0}^{T}d_t
  \{C_{jt}(a)-C_{jt}(0)\}.
\]
If the jurisdiction uses a monetary value or opportunity-cost threshold
$\lambda_j$ per unit of the selected health consequence, incremental net
monetary benefit is \citep{stinnett_mullahy_1998}
\[
\operatorname{INMB}_j(a)=
\lambda_j\Delta Q_j(a)-\Delta C_j(a).
\]
For an eligible population of size $N_{jt}$ and a budget-period cost difference
$\Delta c_{jt}(a)$, a transparent budget-impact calculation is
\[
\operatorname{BI}_j(a)=
\sum_{t=0}^{T_B}N_{jt}\,\Delta c_{jt}(a),
\]
with uptake, displacement and implementation costs modelled explicitly where
relevant. A positive INMB does not imply affordability: an intervention can be
cost-effective over its analytic horizon while imposing a budget impact that a
service cannot absorb. Conversely, a small budget impact does not establish
value.

Distribution cannot be recovered from the population mean. If $g$ indexes
equity-relevant groups, an appraisal should retain at least
\[
\operatorname{Dist}_j(a)=
\left\{\pi_{gj},\Delta Q_{gj}(a),\Delta C_{gj}(a),
\Delta X_{gj}(a)\right\}_{g=1}^{G},
\]
where $\pi_{gj}$ is the population share of group $g$. Distributional
cost-effectiveness analysis may then apply an explicitly chosen inequality
judgement and include the distribution of displaced health; H-Potential does
not supply that judgement automatically. \citep{asaria_2016_dcea}

A small numerical example clarifies the separation. Suppose an intervention
changes the mean profile from $(0.45,0.35,0.30)$ to
$(0.55,0.50,0.40)$, so that
\[
\Delta X_j(a)=(0.10,0.15,0.10).
\]
Suppose also that $\Delta Q_j(a)=0.08$, $\Delta C_j(a)=1{,}200$ monetary
units and $\lambda_j=20{,}000$. Then
\[
\operatorname{INMB}_j(a)
=(20{,}000\times0.08)-1{,}200=400.
\]
At an initial uptake of 10,000 people and an undiscounted first-year incremental
cost of 1,200 per person, the first-year budget impact is 12 million monetary
units. These quantities answer different questions: the profile identifies the
mechanism of change; INMB addresses value relative to opportunity cost; budget
impact addresses fiscal exposure. The numbers are deliberately illustrative,
not estimates from NHATS.

Now consider a second jurisdiction $k$ in which the same profile and health
effects are expected, but $\lambda_k=10{,}000$ and incremental cost is 1,400.
Then
\[
\operatorname{INMB}_k(a)
=(10{,}000\times0.08)-1{,}400=-600.
\]
The intervention has not ceased to work, yet its economic ranking has changed.
The contrast is the practical content of the non-transportability proposition
proved in Appendix~\ref{app:mathematical_foundations}: transport of a clinical
or capability effect does not entail transport of a policy decision. The
remaining coordinates of $\mathcal D_j(a)$ require uncertainty analysis and
deliberative assessment of distribution, legality, accountability and public
reason. They cannot be calculated away.

If the calculation cannot settle the decision, something else must, and that
something is a process rather than a formula.

\section{Policy Bearing and Governance}

The policy attraction of H-Potential is that it can reveal why similar clinical
states lead to different realised opportunities. The corresponding danger is
that a public institution may turn the score into a rank, threshold or automated
eligibility rule. A structurally disadvantaged person could then be described
as having low personal potential and be denied investment precisely because the
system has constrained them.

The correct policy interpretation is remedial rather than desert-based. A low
$P$ may indicate clinical or rehabilitative need. A low $S$ may indicate a need
for support, accessible environments or assistive technology. A low $R$ may
identify institutional barriers requiring policy correction. The profile points
to different duties; it does not measure worth, compliance or probable return
on investment.

The relevant policy problem is global, but it is not context-free. WHO defines
health technology assessment as a multidisciplinary appraisal of clinical,
economic, organisational, social, legal and ethical consequences and urges
Member States to develop methods suited to their national circumstances.
\citep{who_hta_2014,who_hta_institutional_2021} Its recent account of health
equity likewise locates avoidable health differences in the distribution of
power, money, resources and living conditions. \citep{who_sdh_equity_2025}
These propositions support the architecture of H-Potential, in which health
capability is jointly biological and relational. They also defeat the idea of a
single, globally portable decision threshold. A measure can travel as a
hypothesis before it can travel as a calibrated instrument.

This distinction resolves the apparent tension between a general theory and a
United States development study. The conceptual claim is not a claim about one
country: in any society, realised health possibilities may depend on embodied
capacity, enabling relations and institutional arrangements. The empirical
claim is much narrower. NHATS permits an initial test among older Medicare
beneficiaries in the United States. It cannot establish that the same
indicators, loadings, score distributions or prognostic relations hold for
younger people, other disease groups or other health systems. Universality must
therefore be earned through a sequence of transportability tests, not asserted
from the breadth of the theory. \citep{moons_2019,collins_2024}

Four kinds of transport must be distinguished. \emph{Conceptual transport}
asks whether the three-domain account makes sense in the receiving setting.
\emph{Measurement transport} asks whether items have comparable meaning,
requiring qualitative work, measurement-invariance analysis and examination of
differential item functioning. \emph{Predictive transport} asks whether risks
remain calibrated and whether local recalibration or refitting is necessary.
\emph{Decision transport} asks whether an intervention remains worthwhile once
local costs, displaced services, budgets, legal duties and public values are
considered. Success at an earlier level does not entail success at a later one.

Country settings are thus comparative tests rather than rival answers to the
question of scope. The United States supplies the developmental cohort. South
Africa is a substantively important future case because a more resource-
constrained and unequal health system provides a demanding test of the
framework's systemic domain; its emerging HTA arrangements illustrate an
appraisal that keeps effectiveness, equity, cost-effectiveness, budget impact
and societal considerations distinct. \citep{south_africa_hta_2026} Neither
country defines the theory. Each reveals different conditions under which it
might fail, require reinterpretation or acquire policy value.

Accordingly, responsible use cannot be reduced to a procedural checklist. It
requires an evidential chain: participation by affected communities in defining
the domains; external validation and invariance testing; public disclosure of
items, weights, thresholds and missing-data rules; reporting of the domain
profile beside any total; and continuing audit for disparate impact and
self-reinforcing feedback. Where a decision affects an individual, reasons and
an effective route of appeal are indispensable. No score should be the sole
basis for automated denial, rationing or reimbursement. Economic appraisal must
remain separate enough to expose opportunity cost and budget impact rather than
hiding them inside the construct.

These safeguards delimit the policy claim. At the present stage H-Potential is
a research and diagnostic framework. It may inform needs assessment, service
design, inequality surveillance and hypothesis generation. It is not validated
as a benefit-package rule, reimbursement threshold, clinical triage score or
individual entitlement formula.

That leaves the original formula, which needs demoting before it can be
replaced.

\section{The H-potential Equation as a Heuristic, Not an Estimator}

The original formulation
\begin{equation}
H_{\text{potential}}=(L\cdot Q)S-B
\end{equation}
was useful because it expressed an intuitive proposition: biological potential is modified by enabling conditions and constrained by structural burden. The literature, however, does not justify the equation as an empirical law.

$L$ is remaining length of life and $Q$ is the quality of that life, so the
product $L\cdot Q$ is the familiar quality-adjusted survival of health
economics: a clinician can read it as ``how many good years does the biology
allow''. $S$ is a multiplier for social conditions, so a value above one says
that a person's circumstances let them do more with the same body and a value
below one says the reverse. $B$ is a burden imposed by the system, subtracted
at the end.

Multiplication by $S$ assumes that social conditions proportionally scale every unit of biomedical potential. Subtraction of $B$ assumes that systemic burden exists on a common additive scale with the quantity produced by $(L\cdot Q)S$. Neither assumption follows from WHO's framework, capability theory, social-determinants research or frailty science. In particular, capability theory does not imply that a reduction in social opportunity should be numerically equivalent to a unit reduction in biological capacity.

\begin{proposition}[The heuristic is not bounded below by zero]
\label{prop:lit-heuristic}
Let $L,Q,S,B\in[0,1]$ and write $h=(LQ)S-B$. Then
\[
  h\in[-1,1],
\]
and both bounds are attained. In particular $h<0$ whenever $B>(LQ)S$, so $h$
takes negative values on the interior of its domain. If $L$, $Q$, $S$ and $B$
are independent and uniform on $[0,1]$, then
$\Pr(h<0)=1-\mathbb E[LQS]=1-\tfrac18=\tfrac78$. No admissible realisation
function has this property: every aggregator satisfying the normalisation and
capacity-boundedness axioms of Appendix~\ref{app:mathematical_foundations}
returns a value in $[0,P]$.
\end{proposition}

\begin{proof}[Proof sketch]
$(LQ)S$ ranges over $[0,1]$ and $B$ ranges over $[0,1]$ independently, so their
difference ranges over $[-1,1]$; $h=-1$ at $L=Q=S=0$, $B=1$, and $h=1$ at
$L=Q=S=1$, $B=0$. The sign condition is immediate. For the probability, $L$,
$Q$ and $S$ are independent with mean $\tfrac12$, so $\mathbb E[LQS]=\tfrac18$,
and $\Pr(B>LQS)=\mathbb E[1-LQS]$ because $B$ is uniform and independent of the
product.
\end{proof}

\begin{reading}[Why a negative score disqualifies the formula]
A capability of $-0.41$ does not mean anything. There is no state of a person
that is less available to them than nothing, and no clinician can act on the
difference between $-0.41$ and $-0.62$. The subtraction of $B$ therefore fails
before any question of fit arises, because the quantity it produces is not on a
scale a capability could live on.

The seven-eighths figure makes the point sharper. Under an even spread of
inputs the formula is negative most of the time. It is not that the heuristic
occasionally misbehaves at the edges; the well-behaved region is the corner.
That is why the equation is kept as a statement of the intuition and the
empirical model is built from axioms instead.
\end{reading}

\begin{workedexample}[lit-heuristic]{Mrs A scored by the original formula}
Mrs~A's profile is $(P,S,R)=(0.55,0.35,0.40)$, so her structural burden is
$B=1-R=0.60$. Reading her biomedical potential as the product $LQ=0.55$, the
original heuristic gives
\[
  h=(0.55\times0.35)-0.60=0.1925-0.60=-0.4075 .
\]
Her burden would have to fall to $B=0.1925$, meaning $R=0.8075$, before the
heuristic reported a non-negative capability at all. Everything below that is
recorded as a negative number.

The realisation function on the same profile gives
$0.55\sqrt{0.35\times0.40}=0.206$, which sits inside $[0,P]=[0,0.55]$ as the
axioms require, and falls towards zero as the burden rises without ever
crossing it.

\begin{wecode}[Python]
from hpv import crf

heuristic = lambda lq, s, b: lq * s - b
print(f"heuristic {heuristic(0.55, 0.35, 0.60):+.4f}")
print(f"CRF       {crf.H(0.55, 0.35, 0.40):+.4f}")
for b in (0.20, 0.40, 0.60, 0.80):
    print(f"B={b:.2f}  heuristic {heuristic(0.55, 0.35, b):+.4f}  "
          f"CRF {crf.H(0.55, 0.35, 1 - b):+.4f}")
\end{wecode}

\weanswer{$h=-0.4075$ for Mrs~A, against a realisation-function value of
$0.206$. Asserted by
\texttt{test\_proposition\_lit\_heuristic\_admits\_negative\_values} and
\texttt{test\_proposition\_lit\_heuristic\_negative\_on\_most\_of\_the\_cube}.}
\end{workedexample}

The equation should therefore remain as a conceptual heuristic while the empirical model is estimated independently. The thesis should not interpret a standardised SEM factor score as years of healthy life, nor should it convert an arbitrary latent scale into a clinically meaningful unit without external calibration.

Demoting the heuristic leaves a vacancy, which the next section fills.

\section{From heuristic to an axiomatic realisation model}
\label{sec:intro-maths}

The central formal contribution of the thesis is to replace that heuristic with a principled object. Chapter~\ref{chap:crf} develops an axiomatic Capability Realisation Function (CRF): rather than assert a formula, it states the properties a realisation measure must satisfy, including non-compensation, and derives the admissible family from them. The primary form, \citep{cobb_douglas_1928}
\begin{equation}
    H(P,S,R)=P\,S^{\alpha}R^{1-\alpha},\qquad R=1-B,\ 0<\alpha<1,
    \label{eq:crf_intro}
\end{equation}
expresses realised potential $H$ as biomedical potential $P=L\cdot Q$ scaled by a non-compensatory combination of social enablement $S$ and structural freedom $R$.

\begin{reading}[What equation \eqref{eq:crf_intro} says, in words]
Realised health capability is what the body makes possible, discounted by the
fraction of that possibility the person's circumstances actually let them
reach. The discount is the term $S^{\alpha}R^{1-\alpha}$, a weighted geometric
average of social enablement and freedom from structural burden. It runs
between $0$ and $1$, so realised capability can never exceed biomedical
potential: the arithmetic cannot invent health the body does not have. The
exponent $\alpha$ says how the two kinds of circumstance are weighted against
each other, and, more importantly, that they are weighted \emph{multiplicatively}
rather than added.

A geometric average is the formal statement of ``a chain is only as strong as
its weakest link''. Multiply by zero on either side and the product is zero.
This is the property an additive index cannot reproduce, and it is the reason
the form was chosen rather than assumed.
\end{reading}

The three coordinates are used by every chapter that follows, so they are
fixed here once.

\begin{definition}[Domain profile]
\label{def:int:profile}
A person's domain profile is the triple
\[
  \mathbf X=(P,S,R)\in[0,1]^3,
\]
in which $P$ is Biomedical Potential, $S$ is Social Modifiers scored as social
enablement, and $R=1-B$ is structural freedom, the complement of Systemic
Burdens. Every coordinate is oriented so that $0$ is the worst attainable value
and $1$ the best. The realisation function \eqref{eq:crf_intro} is then a map
$H:[0,1]^3\to[0,1]$, and realised capability is bounded by the same interval as
its inputs.

\end{definition}

\begin{remark}[Two scales for $P$, and the conversion between them]
\label{rem:int:scales}
Biomedical potential is reported on two scales in this thesis, and the reader
should know which is in use before meeting either.

The \emph{unit scale} of Definition~\ref{def:int:profile} puts $P\in[0,1]$,
normalised within the calibration population. Appendix~\ref{app:mathematical_foundations}
and the empirical chapters use it, because a latent measurement model returns a
relative position and not a quantity of anything.

The \emph{natural scale} puts $P=L\cdot Q$ in quality-adjusted life-years, as
Chapter~\ref{app:crf} does when the realisation function has to meet an
economic evaluation. Here $H$ inherits the units of $P$ and is read in good
years.

The two differ by normalisation alone. Because the enablers enter only through
the bounded multiplier $S^{\alpha}R^{1-\alpha}$, which is invariant to the
scale of $P$, any statement about the \emph{ratio} of two people's realised
capability, and the whole of the gap decomposition, is identical on both. Only
the level changes.

Mrs~A is the worked case for this. On the natural scale, with $P=8.4$ good
years, she realises $3.1$ of them. On the unit scale, with $P=0.55$, she scores
$0.206$. These are the same calculation and the same woman;
$0.206$ is not a number of years and $3.1$ is not a position in a
distribution. Where a figure for her appears later without its scale named,
$0.206$ is the unit-scale value.

Carrying two scales is a cost, not a feature. It is paid because the
measurement model cannot deliver years without an external calibration study
that has not been done, while an economic evaluation cannot use a position in a
distribution. \S\ref{sec:intro-maths} returns to why the thesis refuses to
manufacture the missing calibration.
\end{remark}

The choice matters clinically, and the two running cases show why. Take Mrs~A
first, with biomedical potential $P=8.4$ good years as above. Score her social
enablement at $S=0.35$ (she lives alone, her daughter is three hours away, she
has withdrawn from the church group she used to attend) and her structural
freedom at $R=0.40$ (no car, two buses, missed appointments, a co-payment she
notices). With $\alpha=0.5$ the multiplier is
$\sqrt{0.35\times0.40}=0.37$, and her realised capability is
$8.4\times0.37=3.1$ good years. An additive index that averaged her
circumstances would give $(0.35+0.40)/2=0.375$ and $3.2$ good years. The two
agree, which is exactly the uninformative case: her circumstances are equally
poor on both sides, so no rule for trading them can matter.

Now change one thing. Suppose transport is solved: a community minibus takes
her to clinic, so $R$ rises from $0.40$ to $0.90$, while her isolation is
untouched at $S=0.35$. The geometric form gives
$\sqrt{0.35\times0.90}=0.56$ and realised capability $4.7$ good years; the
additive form gives $0.625$ and $5.3$. The gap has opened, and it opens further
as one term approaches zero. Suppose instead that Mrs~A's daughter dies and she
is left with no informal support at all, $S=0.05$, while the minibus keeps
$R=0.90$. The geometric form returns $\sqrt{0.05\times0.90}=0.21$, so $1.8$
good years. The additive form returns $0.475$, so $4.0$ good years, and would
record her as more than twice as well off as she is. The additive index has
allowed excellent transport to compensate for having nobody. That substitution
is precisely what the non-compensation axiom forbids, and it is why the
functional form is derived rather than picked.

\begin{workedexample}[int:mrsa]{Mrs~A under the geometric rule and the additive rule}
Take $P=8.4$ good years and $\alpha=0.5$, and compare three states of Mrs~A's
circumstances. The question is how far the two rules disagree, and where.

As she is, $S=0.35$ and $R=0.40$. The geometric multiplier is
$\sqrt{0.35\times0.40}=0.374$ and the additive one is $(0.35+0.40)/2=0.375$, so
the two rules return $3.14$ and $3.15$ good years. Neither rule is doing any
work here, because there is nothing to trade.

Give her the community minibus, $R=0.90$. The multipliers separate to $0.561$
and $0.625$, which is $4.71$ good years against $5.25$. Now take the support
away instead, $S=0.05$ with the minibus kept at $R=0.90$. The geometric
multiplier falls to $0.212$ and the additive one only to $0.475$: $1.78$ good
years against $3.99$. The additive rule records her as $2.24$ times as well off
as the geometric rule does, having let transport stand in for having nobody.

\begin{wecode}[Python]
from hpv.crf import contextual_mean, H

P = 8.4
for s, r in ((0.35, 0.40), (0.35, 0.90), (0.05, 0.90)):
    geo = contextual_mean(s, r, rho=0.0)   # the CRF, a geometric mean
    add = contextual_mean(s, r, rho=1.0)   # the additive comparator
    print(f"S={s:.2f} R={r:.2f}: CRF {P * geo:.2f} yr   "
          f"additive {P * add:.2f} yr   overstatement {add / geo:.2f}x")

print(H(0.55, 0.35, 0.40))   # the same profile on the unit scale: 0.2058
\end{wecode}

\weanswer{$3.14$, $4.71$ and $1.78$ good years under the CRF, against $3.15$,
$5.25$ and $3.99$ under the additive rule. At $S=0.05$ the additive rule
overstates Mrs~A by a factor of $2.24$. Asserted by
\texttt{test\_worked\_example\_intro\_mrs\_a\_geometric\_versus\_additive}.}
\end{workedexample}

Mr~B shows the same structure from the other end. His biomedical potential is
poor and well characterised: an inhibitor, a destroyed knee, chronic pain.
Suppose $P=4.0$. In a system that funds emicizumab and staffs a haemophilia
comprehensive-care centre, his structural freedom is high, $R=0.85$, and stable
employment and a partner give $S=0.70$; his realised capability is
$4.0\sqrt{0.70\times0.85}=3.1$, close to the ceiling his biology allows. Move
him to a jurisdiction where bypassing agents are rationed and the nearest centre
is in another province, $R=0.20$, and the same body yields
$4.0\sqrt{0.70\times0.20}=1.5$. Nothing about Mr~B has changed. The formula
attributes the entire difference to the system, which is the claim the thesis
is making and which \S\ref{sec:intro-maths} must therefore be able to defend
rather than assert. This single object does the work the heuristic could not: it proves the biomedical--context complementarity that is the mathematical form of the intersectional \citep{crenshaw_1991,bowleg_2008} claim, decomposes any capability gap exactly into biomedical, social and structural parts, and, through an equity-sensitive population functional and a budget-constrained policy gradient, turns individual capability into an integrated, distribution-aware policy calculation. It is also falsifiable: the chapter pre-specifies the head-to-head tests that would reject or narrow it. The construct's meaning and the CRF's parameters are fixed without reference to the outcomes, so the empirical chapters test the model rather than fit it.

\begin{workedexample}[int:mrb]{Mr~B funded and Mr~B rationed, one coordinate apart}
On the unit scale Mr~B's profile is $(0.30,0.70,0.85)$ where his treatment is
funded and $(0.30,0.70,0.20)$ where it is rationed. His biology and his
household are identical in the two; only $R$ moves.

The CRF returns $H=0.231$ funded and $H=0.112$ rationed, a ratio of $2.06$. In
the good-years units used above, with $P=4.0$, the same two multipliers give
$3.09$ years and $1.50$ years. The exact gap decomposition of
Chapter~\ref{chap:crf} splits the log ratio
$\log(0.231/0.112)=0.723$ into three parts. The biomedical part is
$\log(0.30/0.30)=0$ and the social part is $0.5\log(0.70/0.70)=0$. The
structural part, $0.5\log(0.85/0.20)=0.723$, is the whole of it.

\begin{wecode}[Python]
from hpv.crf import H, gap_decomposition

funded, rationed = (0.30, 0.70, 0.85), (0.30, 0.70, 0.20)
print(H(*funded), H(*rationed), H(*funded) / H(*rationed))

bio, soc, structural, total = gap_decomposition(funded, rationed)
print(f"biomedical {bio:.3f}  social {soc:.3f}  "
      f"structural {structural:.3f}  total {total:.3f}")
\end{wecode}

\weanswer{$H=0.231$ funded and $0.112$ rationed, a ratio of $2.06$. The
decomposition puts $0.723$ of a total of $0.723$ on the health system and
nothing on Mr~B. Asserted by
\texttt{test\_worked\_example\_intro\_mr\_b\_funded\_versus\_rationed}.}
\end{workedexample}

This distinction also resolves a potentially serious interpretive problem. A standardised latent factor is not automatically measured in years of healthy life. Unless an external calibration study establishes a mapping between the factor and a time-based health-expectancy scale, $H_{\text{potential}}$ must be reported as a latent score with an explicitly defined direction and uncertainty. The thesis will not infer a ``number of healthy years'' from an arbitrary standardisation.

With the formal object in place, the candidate representations that compete with it, and the tests that would choose between them, can be set out.

\section{Candidate Mathematical Representations and Their Tests}
\label{sec:lit-candidates}

Equation~\ref{eq:hp_general_class} defines a class of admissible models rather
than selecting one. Four candidates expose the principal aggregation assumptions \citep{arrow_chenery_1961}:
\begin{align}
H_{P,i}^{(A)}&=\beta_0+\beta_1BP_i+\beta_2SM_i-\beta_3SB_i+\varepsilon_i,
\label{eq:hp_additive}\\
H_{P,i}^{(I)}&=\beta_0+\beta_1BP_i+\beta_2SM_i-\beta_3SB_i
+\beta_4(BP_iSM_i)+\varepsilon_i,
\label{eq:hp_interaction}\\
H_{P,i}^{(M)}&=BP_i^{\alpha}SM_i^{\beta}(1-SB_i)^{\gamma},
\qquad \alpha,\beta,\gamma>0,
\label{eq:hp_multiplicative}\\
H_{P,i}^{(C)}&=BP_i\left[\delta SM_i^{\rho}
+(1-\delta)(1-SB_i)^{\rho}\right]^{1/\rho},
\quad 0<\delta<1.
\label{eq:hp_ces}
\end{align}
The bounded complement $1-SB_i$ makes zero a complete contextual bottleneck.
All components must be directionally harmonised and scaled before comparison.

The additive model assumes constant marginal contributions and complete
compensation on a common scale. The interaction model tests whether enablement
changes the contribution of biological capacity; $\beta_4=0$ is the nested null
of no complementarity. The multiplicative model imposes positive monotonicity
and proportional trade-offs. The CES \citep{arrow_chenery_1961} model varies substitutability:
$\rho\rightarrow1$ approaches a weighted arithmetic mean,
$\rho\rightarrow0$ the geometric case, and $\rho\rightarrow-\infty$ a
bottleneck limit. These are empirical and normative alternatives.

\begin{proposition}[Aggregation can reverse a ranking]
\label{prop:lit-reversal}
There are profiles that an additive rule and a bottleneck rule order in
opposite directions. Let $x=(1,1,0.2)$ and $y=(1,0.55,0.55)$. The equal-weight
arithmetic contextual mean gives $0.60$ for $x$ and $0.55$ for $y$, so $x$ is
ranked first. The bottleneck rule gives $0.20$ for $x$ and $0.55$ for $y$, so
$y$ is ranked first. The choice of $\rho$ in \eqref{eq:hp_ces} therefore
determines the ordering of two people and is not a technical detail. The proof
is in Appendix~\ref{app:mathematical_foundations}.
\end{proposition}

\begin{reading}[Who the two profiles are]
Profile $x$ is someone with excellent support and a system that fails them
almost completely. Profile $y$ is someone with middling amounts of both. An
additive rule says the first is better placed, because the support makes up for
the failure. A bottleneck rule says the second is, because nothing makes up for
a system score of $0.2$.

Both statements are coherent and they cannot both guide the same service. The
question of which is right is ethical, and no amount of data settles it. What
the data can do is show how far the answer moves as $\rho$ moves, which is why
substitution sensitivity is reported alongside every total in this thesis.
\end{reading}

\begin{workedexample}[lit-reversal]{The reversal, with the numbers}
Take $x=(1,1,0.2)$ and $y=(1,0.55,0.55)$, both with full biomedical capacity so
that the contextual mean is the whole story.

The arithmetic mean of the two context scores is $\tfrac12(1+0.2)=0.60$ for $x$
and $\tfrac12(0.55+0.55)=0.55$ for $y$. The geometric mean is
$\sqrt{1\times0.2}=0.447$ for $x$ and $\sqrt{0.55\times0.55}=0.55$ for $y$. The
bottleneck value is $\min(1,0.2)=0.20$ for $x$ and $0.55$ for $y$.

The additive rule puts $x$ ahead by $0.05$. Every rule with $\rho\le0$ puts $y$
ahead, by $0.10$ in the geometric case and by $0.35$ at the bottleneck limit.
The gap between the two verdicts is larger than either margin.

\begin{wecode}[Python]
from hpv import crf

x, y = (1.0, 1.0, 0.2), (1.0, 0.55, 0.55)
for name, rho in (("arithmetic", 1.0), ("geometric", 0.0), ("bottleneck", -800.0)):
    hx, hy = crf.H(*x, rho=rho), crf.H(*y, rho=rho)
    print(f"{name:>11}: x {hx:.3f}  y {hy:.3f}  winner {'x' if hx > hy else 'y'}")
\end{wecode}

\weanswer{$x$ scores $0.60$, $0.447$ and $0.20$ as $\rho$ falls from $1$ to the
bottleneck limit, while $y$ scores $0.55$ throughout: the ranking reverses
between the first rule and the other two. Asserted by
\texttt{test\_worked\_example\_lit\_ranking\_reversal}.}
\end{workedexample}

For the heuristic $(LQ)S-B$,
\begin{equation}
\frac{\partial H_P}{\partial L}=QS,
\qquad
\frac{\partial^2 H_P}{\partial L\,\partial S}=Q.
\label{eq:heuristic_derivatives}
\end{equation}
The positive cross-partial is a substantive complementarity claim: support
raises the marginal contribution of biological capacity. Additive barriers also
require commensurability between $B$ and $(LQ)S$. Failure of either assumption
rejects that functional form, not automatically the general construct.

The transition from construct to measurement is separate. For a reflective
domain $\eta_i$, indicator $j$ is represented as \citep{bollen_1989,kline_2015}
\begin{equation}
x_{ij}=\nu_j+\lambda_j\eta_i+\epsilon_{ij},
\label{eq:reflective_measurement_bridge}
\end{equation}
whereas a formative composite reverses the direction,
$\eta_i=\sum_j w_jx_{ij}+\zeta_i$. The code therefore does not force social
and structural indicators into the reflective model. It compares admissible
measurement structures without outcomes and only then compares the candidate
realisations against prespecified comparators and downstream outcomes.

The more defensible conceptual representation is therefore
\begin{equation}
H_{\text{potential}}=f(BP,SM,SB),
\end{equation}
where $BP$ denotes Biomedical Potential, $SM$ Social Modifiers and $SB$ Systemic Burdens, and $f$ is intentionally left unspecified at the theoretical stage. The empirical analysis is then responsible for determining whether these domains share a common higher-order structure and whether a scalar representation is warranted.

Everything in this chapter now reduces to a short list of statements that could
turn out to be false.

\section{From Literature to Testable Propositions}

The literature synthesis generates the following propositions.

\begin{enumerate}
    \item \textbf{Multidimensionality proposition.} The indicators selected for the proposed construct will represent more than one empirically distinguishable domain of health-related capability rather than a single undifferentiated health state.
    \item \textbf{Capacity proposition.} Biomedical Potential will show substantial convergence with established measures of intrinsic capacity and current functioning while retaining information not reducible to either.
    \item \textbf{Context proposition.} Social Modifiers and Systemic Burdens will demonstrate associations with health-related functioning that are not fully explained by biomedical capacity alone.
    \item \textbf{Distinctiveness proposition.} The proposed construct will demonstrate incremental information relative to pre-specified frailty and social-vulnerability comparators; otherwise its justification as a distinct metric will be limited.
    \item \textbf{General-factor proposition.} A higher-order H-potential factor will be retained only if the data demonstrate sufficient commonality, factor determinacy and substantive coherence. A correlated-domain model remains a valid endpoint if those conditions are not met.
    \item \textbf{Equity proposition.} The measurement structure will be sufficiently invariant across major demographic and socioeconomic groups to support meaningful comparisons. Where invariance fails, the thesis will report the failure rather than treating observed group differences as direct evidence of unequal health potential.
    \item \textbf{Prognostic proposition.} If a general H-potential score is empirically defensible, it will provide prognostic information for mortality and/or hospitalisation beyond established health and vulnerability measures. This proposition is secondary to construct validity and will not be used to select the measurement model.
\end{enumerate}

Each proposition has a corresponding empirical obligation. The first requires exploratory and confirmatory dimensionality analysis; the second requires convergent and discriminant validity testing; the third requires domain-level and sensitivity analyses; the fourth requires nested and incremental prognostic comparisons; the fifth requires competing measurement structures and factor-determinacy assessment; the sixth requires measurement-invariance analysis; and the seventh requires independently specified survival and hospitalisation models with internal validation.

This mapping also establishes a direct link to the computational implementation. The code should not merely calculate a final score. It must preserve the distinction between theory generation, measurement evaluation, score estimation and outcome prediction. Every model-selection decision must therefore be traceable to a stated theoretical proposition or methodological criterion.

Traceability is a property of the code as much as of the argument, which makes
it a design requirement rather than a promise.

\section{Implications for the Empirical Design and Codebase}

The literature review changes the empirical design in several concrete ways. First, indicator selection must be documented as a theory-to-variable mapping. Each indicator should have a stated construct, domain, expected direction and justification from the literature. Second, candidate measurement models must be fitted before outcome variables are introduced into model selection. Third, the analysis must include explicit tests of alternative structures rather than only the preferred model.

Fourth, the code must interrogate the possibility that different domains have different measurement ontologies. Biomedical indicators may plausibly be reflective manifestations of underlying capacity, whereas social conditions may function as formative contributors, contextual variables or external causes. The analysis should therefore avoid forcing all indicators into one reflective CFA merely because a common SEM framework is available.

Fifth, the analysis must distinguish model fit from substantive validity. A model can fit well and still measure the wrong construct. Conversely, a theoretically coherent construct may show imperfect fit because the observed indicators are necessarily heterogeneous. Fit indices should therefore be interpreted alongside reliability, factor determinacy, construct validity, invariance and comparison with established instruments.

Sixth, the prognostic analysis must remain downstream of measurement. Mortality and hospitalisation should not be used to discover the latent structure. They should test whether the construct developed independently of the outcomes contains useful information. Comparisons with frailty, intrinsic capacity and social vulnerability are therefore not optional sensitivity analyses; they are central tests of the thesis's claim to added value.

Finally, the code should be designed to expose failure. It should be possible for the pipeline to conclude that no defensible general factor exists, that a domain is empirically redundant, that measurement invariance fails, or that prognostic performance adds little beyond established comparators. Such outcomes would not constitute programming failures. They would be scientifically informative findings about the proposed construct.

\section{A reader's glossary: the same idea in two languages}
\label{sec:intro-glossary}

This thesis is written for three readerships that do not share a vocabulary.
Clinicians, health economists and statisticians each have a name for most of
the ideas used here, and the names are different. Table~\ref{tab:bilingual}
pairs them, so that a term met later in an unfamiliar form can be looked up
rather than guessed at. Each pairing is introduced again in the text at the
point where it first matters; the table is for reference, not instruction.

The pairings are not exact synonyms. They are close enough that a reader who
understands one side will not be lost on the other, which is all a glossary can
honestly do.

\begin{table}[htbp]
\centering
\caption{Quantitative terms and their clinical counterparts, as used in this thesis.}
\label{tab:bilingual}
\footnotesize
\renewcommand{\arraystretch}{1.25}
\begin{tabular}{p{3.3cm} p{3.9cm} p{6.2cm}}
\toprule
\textbf{Quantitative term} & \textbf{Clinical counterpart} & \textbf{What both are naming} \\
\midrule
Endogeneity & Confounding by indication & The exposure was chosen because of something that also affects the outcome. \\
\addlinespace
Attenuation bias & Regression dilution & Measuring the predictor with error pulls its estimated effect towards no effect. \\
\addlinespace
Shrinkage & Regression to the mean & An extreme estimate from sparse information is pulled back towards the average. \\
\addlinespace
Censoring & Alive at last contact & The person's later fate is unknown, which is information, not a missing value. \\
\addlinespace
Competing risk & An event that ends the follow-up & Death removes a person from being at risk of the admission you were counting. \\
\addlinespace
Immortal time & A window where the event could not occur & The design guarantees survival over an interval, and the guarantee looks like benefit. \\
\addlinespace
Discrimination & Ranking patients correctly by risk & Whether the sicker of any two patients gets the higher score. \\
\addlinespace
Calibration & Whether a stated 10 per cent means one in ten & Whether the predicted probabilities are the observed frequencies. \\
\addlinespace
Incremental value & Does this change management? & Whether the new measure adds anything to the one already in use. \\
\addlinespace
Discriminant validity & Is this troponin measured a second way? & Whether two instruments are the same thing under different names. \\
\addlinespace
Reflective indicators & Symptoms of one underlying disease & Tiredness, cold intolerance and dry skin as effects of a low thyroid. \\
\addlinespace
Formative indicators & Ingredients of a risk score & Age, blood pressure, lipids and smoking, which need not correlate. \\
\addlinespace
Latent variable & A quantity inferred, never measured & Glomerular filtration rate behind creatinine and cystatin~C. \\
\addlinespace
Identification & Telling a blocked patient from a choosing one & Whether the data can distinguish the explanations of the same behaviour. \\
\addlinespace
Non-compensation & A chain is as strong as its weakest link & Cardiac output cannot rescue a haemoglobin of two. \\
\addlinespace
Elasticity of substitution & How far one thing stands in for another & Whether good support can make up for a system that will not treat you. \\
\addlinespace
Opportunity cost, $\lambda$ & Patients you never meet & The health lost elsewhere when this budget is spent here. \\
\addlinespace
Inequality aversion, $\varepsilon$ & How much a gain to the worst-off counts & The exchange rate between total health and its distribution. \\
\addlinespace
Net benefit & Finding a case against over-treating & The trade a clinician makes at a chosen decision threshold. \\
\addlinespace
Internal validation & The model marking its own homework & How much of the apparent performance will survive a new sample. \\
\bottomrule
\end{tabular}
\end{table}

\section{Conclusion}

The literature reveals a mature field rather than an empty space. QALYs and DALYs address health loss and comparative valuation; health expectancy addresses expected population health over time; quality-of-life instruments describe experienced health; frailty captures accumulated vulnerability; intrinsic capacity describes the physical and mental capacities available to an individual; capability measures address opportunities to achieve valued states; and social-vulnerability measures quantify contextual disadvantage. The challenge is not to replace these traditions but to determine whether a coherent bridge among them is empirically defensible.

The proposed H-potential therefore enters the literature with a deliberately modest claim. Its novelty, if established, will not be that it is the first multidimensional health metric. Nor will it be that it includes social determinants or uses latent-variable methods. The potential contribution is the explicit testing of a health-bounded construct in which biological capacity, enabling social conditions and structural constraints are considered jointly, while measurement validity is kept analytically separate from prognostic performance.

That proposition remains open. If the data show that H-potential is indistinguishable from intrinsic capacity, frailty, social vulnerability or an established health metric, the thesis must conclude that the proposed integration does not warrant a separate construct. If the domains are coherent but a general factor is not, the appropriate result may be a multidomain profile rather than a scalar score. If the score is coherent and incrementally predictive, the contribution becomes stronger. In all cases, the scientific value lies in allowing the empirical evidence to determine which interpretation survives interrogation.
